\DocumentMetadata{
  pdfversion=2.0,
  pdfstandard=ua-2,
  lang=en-US,
  tagging=on,
  tagging-setup={math/setup=mathml-SE,tabsorder=structure}
}
\documentclass[11pt,letterpaper]{article}
\usepackage[margin=0.68in,includefoot]{geometry}
\usepackage{newcomputermodern}
\usepackage{amsmath,amscd,mathtools}
\providecommand{\square}{\mdlgwhtsquare}
\usepackage[hidelinks]{hyperref}
\hypersetup{
  pdftitle={Logic PhD Exam, August 2022},
  pdfauthor={University of Florida Department of Mathematics},
  pdfsubject={Graduate mathematics examination},
  pdfdisplaydoctitle=true,
  bookmarksopen=true
}
\setcounter{secnumdepth}{0}
\setlength{\parindent}{0pt}
\setlength{\parskip}{0.28em}
\setlength{\emergencystretch}{3em}
\setlength{\leftmargini}{2.15em}
\setlength{\leftmarginii}{2.65em}
\setlength{\leftmarginiii}{3em}
\renewcommand{\labelenumi}{\textbf{\theenumi.}}
\pagestyle{empty}
\raggedbottom
\allowdisplaybreaks
\newenvironment{examproblems}[1]{%
  \begin{enumerate}%
  \setcounter{enumi}{#1}%
  \setlength{\topsep}{0.35em}%
  \setlength{\partopsep}{0pt}%
  \setlength{\itemsep}{0.48em}%
  \setlength{\parsep}{0.18em}%
}{\end{enumerate}}
\newenvironment{examsubparts}{%
  \begin{enumerate}%
  \setlength{\topsep}{0.18em}%
  \setlength{\partopsep}{0pt}%
  \setlength{\itemsep}{0.16em}%
  \setlength{\parsep}{0.1em}%
}{\end{enumerate}}
\newenvironment{examinstructions}{%
  \par\begingroup\itshape
}{%
  \par\endgroup\vspace{0.18em}
}
\makeatletter
\renewcommand\section{\@startsection{section}{1}{0pt}%
  {-1.25ex plus -.25ex minus -.1ex}%
  {0.55ex plus .1ex}%
  {\normalfont\large\bfseries}}
\renewcommand{\@maketitle}{%
  \newpage\null\vskip -1.2em%
  {\footnotesize\bfseries
    \href{https://www.ufl.edu/}{University of Florida}, \href{https://math.ufl.edu/}{Department of Mathematics}\par}%
  \vskip 0.18em%
  \tagstructbegin{tag=Title}%
    {\LARGE\bfseries\href{https://gma.math.ufl.edu/exams/logic-phd/}{Logic PhD Exam}, August 2022\par}%
  \tagstructend%
  \vskip 0.35em%
  \tagmcbegin{artifact}%
    \hrule height 1pt%
  \tagmcend%
  \vskip 0.55em%
}
\makeatother
\title{Logic PhD Exam, August 2022}
\author{}
\date{}
\begin{document}
\maketitle
\thispagestyle{empty}
\begin{examinstructions}
Solve 5 of the following problems; at least one from each section.
\end{examinstructions}
\section{A. Set Theory.}
\begin{examproblems}{0}
\item
Let~\(X\) be an uncountable set and let \(F:X\to X\) be a function. Show that there are distinct elements \(x,y\in X\) such that \(x\ne F(y)\) and \(y\ne F(x)\).
\item
Let~\(P\) be a forcing. Show that there is an uncountable cardinal~\(\kappa\) such that \(P\Vdash \check{\kappa}\) is a cardinal.
\item
Provide an example of a subset of a Polish space which is analytic and not Borel. Provide a proof of these two properties of your set.
\end{examproblems}
\section{B. Computability.}
\begin{examproblems}{0}
\item
Give an example of a subset of \(\mathbb N\) that is \(\Pi^0_1\) and not \(\Sigma^0_1\) (and prove that your set has these two properties).
\item
Prove that there exist incomparable Turing degrees.
\item
Prove that there exist disjoint computably enumerable sets \(A,B\subseteq\mathbb N\) such that there is no computable set~\(C\) which contains~\(A\) and is disjoint from~\(B\).
\end{examproblems}
\section{C. Model Theory.}
\begin{examproblems}{0}
\item
Let \(\mathcal L\) be a language comprised of a single binary relation symbol. Give an example of a countably infinite \(\mathcal L\)-structure \(\mathcal M\) such that. Make sure to prove both statements (i) and (ii).
\begin{examsubparts}
\item[\textnormal{(i)}]
\(\operatorname{Th}(\mathcal M)\) is \(\aleph_0\)-categorical, and
\item[\textnormal{(ii)}]
\(\operatorname{Th}(\mathcal M)\) does not have quantifier elimination.
\end{examsubparts}
\item
Let \(\mathcal L\) be a language comprised of two unary relation symbols~\(G\) and~\(B\), and a binary relation symbol~\(R\). Let \(\mathcal K\) be the class of all finite graphs with a bipartition, i.e., \(\mathcal K\) consists of all finite \(\mathcal L\)-structures \(\mathcal M=(M,G^{\mathcal M},B^{\mathcal M},R^{\mathcal M})\) such that. Prove or disprove: \(\mathcal K\) is a Fraïssé class.
\begin{examsubparts}
\item[\textnormal{(i)}]
\(R^{\mathcal M}\) is a symmetric and irreflexive relation on~\(M\),
\item[\textnormal{(ii)}]
\(G^{\mathcal M}\) and \(B^{\mathcal M}\) form a partition of~\(M\),
\item[\textnormal{(iii)}]
if \((x,y)\in R^{\mathcal M}\) then exactly one of \(x,y\) belongs to \(G^{\mathcal M}\) and the other belongs to \(B^{\mathcal M}\).
\end{examsubparts}
\item
Let \(\mathcal L=\{R\}\) be the language of graphs (i.e., \(R\) is a binary relation symbol). Let \(\mathcal K_0\) be the class of all connected graphs, and let \(\mathcal K_1\) be the class of all disconnected graphs. Prove or disprove each of the following:
\begin{examsubparts}
\item[\textnormal{(i)}]
\(\mathcal K_0\) is an elementary class;
\item[\textnormal{(ii)}]
\(\mathcal K_1\) is an elementary class.
\end{examsubparts}
\end{examproblems}
\end{document}
